\begin{frame}{Наблюдательный тракт ECAP-like}
\begin{itemize}
  \item потенциалы действия и задержки проведения;
  \item пространственное суммирование и геометрия электродов;
  \item объёмная проводимость;
  \item стимуляционный артефакт, физиологический и аппаратный шум;
  \item насыщение и пропуски.
\end{itemize}
До количественной проверки используется только термин \textbf{ECAP-like}.
\Evidence{S105, S159, S162}
\end{frame}
\begin{frame}{Четыре самостоятельных эксперимента}
\begin{enumerate}
  \item Детекция параметров воздействия в симуляции.
  \item Оценка нейронной или защитной реакции.
  \item Перенос на одну человеческую целевую переменную.
  \item Условный прогноз ответа на SCS на новых пациентах.
\end{enumerate}
\end{frame}
\begin{frame}{STOP-критерии}
\small
\begin{itemize}
  \item \texttt{STOP-H1}: невоспроизводимость или leakage.
  \item \texttt{STOP-H2-A}: нет обоснованного общего пространства.
  \item \texttt{STOP-H2-B}: нет прироста к human-only и контролям.
  \item \texttt{STOP-H3-A}: нет patient-linked SCS/ECAP и исходов.
  \item \texttt{STOP-H3-B}: эффект исчезает при patient holdout.
\end{itemize}
\end{frame}
